\section{Introduction}
\label{sec:intro}

Tool calling is the primary interface through which large language models (LLMs) act on external environments, and its form has been evolving.
Early agents issue a single tool call per model turn and read its result before deciding the next call, following the reasoning-and-acting paradigm~\citep{yao2023react,schick2023toolformer,qin2024toolllm}.
Recent systems instead adopt \emph{Programmatic Tool Calling} (PTC), where the model writes a program that composes many tool calls with loops, conditionals, and intermediate data processing, and the program runs in a persistent runtime whose variables carry across turns~\citep{gao2023pal,wang2024codeact,smolagents2025,anthropic2025ptc}.
By packing a coherent phase of work into one executable block, PTC reduces the number of model-tool interaction rounds, expresses control flow that a call-per-turn protocol cannot, and lets a later block reuse variables produced by an earlier one.
However, the expressiveness that PTC gains inside a block is not matched by how it organizes work across blocks.

Across blocks, the history that a PTC agent conditions on is a flat sequence of code, standard output, and error text ordered by time, which is a lossy record of what execution did.
A single block on the AppWorld benchmark~\citep{trivedi2024appworld} issues five to ten tool calls whose read or write effects and produced values are flattened into one output string, and a task spans more than one hundred such values across a dozen blocks.
When the agent later needs to know whether a write took effect, which earlier value a new step depends on, or whether two differently phrased queries returned the same content, it has to reconstruct these relations by re-reading raw text.
This difficulty stems from a representation that stores execution as concatenated strings and never as the dependency structure that execution induces.

\begin{figure}[t]
\centering
\begin{tikzpicture}[
  font=\scriptsize,
  block/.style={rounded corners=2pt, draw=black!70, fill=blue!6, minimum width=1.05cm, minimum height=0.45cm, align=center},
  gnode/.style={circle, draw=black!70, minimum size=0.34cm, inner sep=0.5pt, align=center},
  task/.style={rounded corners=2pt, draw=black!80, fill=black!8, minimum width=0.95cm, minimum height=0.4cm, align=center},
  art/.style={gnode, fill=green!12},
  call/.style={gnode, fill=orange!16},
  st/.style={gnode, fill=purple!12},
  fail/.style={gnode, fill=red!16},
  intent/.style={rounded corners=2pt, draw=black!70, fill=yellow!18, minimum width=0.9cm, minimum height=0.38cm, align=center},
  rel/.style={->, >=Stealth, draw=black!55, thin},
  seq/.style={->, >=Stealth, draw=black!75},
]

% ============ Left panel: PTC linear history ============
\node[align=center, font=\scriptsize\bfseries] at (1.7,2.55) {PTC: Linear Execution History};

\node[block] (b1) at (0.55,1.7) {Block $P_1$};
\node[block] (b2) at (0.55,0.9) {Block $P_2$};
\node[block] (b3) at (0.55,0.1) {Block $P_3$};
\node[draw=black!30, fill=black!3, rounded corners=1pt, minimum width=1.35cm, minimum height=0.4cm, align=center] (o1) at (2.35,1.7) {\tiny stdout+err};
\node[draw=black!30, fill=black!3, rounded corners=1pt, minimum width=1.35cm, minimum height=0.4cm, align=center] (o2) at (2.35,0.9) {\tiny stdout+err};
\node[draw=black!30, fill=black!3, rounded corners=1pt, minimum width=1.35cm, minimum height=0.4cm, align=center] (o3) at (2.35,0.1) {\tiny stdout+err};
\draw[seq] (b1)--(o1); \draw[seq] (b2)--(o2); \draw[seq] (b3)--(o3);
\draw[seq] (o1) to[out=-90,in=90] (b2);
\draw[seq] (o2) to[out=-90,in=90] (b3);
\node[align=center, text=red!60!black, font=\tiny] at (1.7,-0.6) {cross-block relations flattened into text,\\ progress and dependencies implicit};

% divider
\draw[black!25, dashed] (3.5,-0.85) -- (3.5,2.75);

% ============ Right panel: GraphPTC ============
\node[align=center, font=\scriptsize\bfseries] at (8.05,2.55) {\method: Dynamic Execution Graph};

\node[task] (tk) at (8.05,1.95) {task};
\node[intent] (in) at (5.5,1.95) {intent\\ \tiny (a,z,e)};

\node[block] (gb1) at (5.5,1.15) {Block $P_1$};
\node[block] (gb2) at (5.5,0.15) {Block $P_2$};

\node[call] (c1) at (6.7,1.5) {\tiny r};
\node[call] (c2) at (6.7,0.85) {\tiny w};
\node[art]  (a1) at (7.7,1.5) {\tiny $v_1$};
\node[st]   (s1) at (7.7,0.85) {\tiny st};
\node[call] (c3) at (6.7,0.15) {\tiny r};
\node[art]  (a2) at (7.7,0.15) {\tiny $v_2$};
\node[fail] (f1) at (6.0,-0.55) {\tiny F};

% intent binding
\draw[rel] (in)--(gb1) node[midway,left,font=\tiny]{impl.};
\draw[rel, dashed] (in) to[out=0,in=180] (tk);

% block1 internals
\draw[rel] (gb1)--(c1) node[midway,above,font=\tiny]{exec};
\draw[rel] (gb1)--(c2);
\draw[rel] (c1)--(a1) node[midway,above,font=\tiny]{prod};
\draw[rel] (c2)--(s1) node[midway,above,font=\tiny]{mut};

% block2 internals + cross-block
\draw[rel] (gb2)--(c3);
\draw[rel] (c3)--(a2);
\draw[rel] (a1) to[out=-90,in=90] (c3) node[midway,right,font=\tiny]{consumes};
\draw[rel, dotted] (a2) to[out=200,in=-20] (a1) node[midway,below,font=\tiny]{equiv.};
\draw[rel] (gb2)--(f1) node[midway,left,font=\tiny]{fails\_at};

% blocks belong to task
\draw[rel, black!35] (tk) to[out=-120,in=60] (gb1);
\draw[rel, black!35] (tk) to[out=-110,in=70] (gb2);

% PATCH loop back
\draw[->, >=Stealth, red!65!black, thick] (f1) to[out=180,in=-140] node[midway,below,font=\tiny,text=red!65!black]{PATCH} (gb2.west);

\node[align=center, text=green!45!black, font=\tiny] at (8.05,-0.7) {typed in-block effects + multi-relation\\ cross-block graph, addressable and revisitable};

\end{tikzpicture}
\vspace{-2mm}
\caption{PTC records execution as a linear stream of code and output where cross-block dependencies and progress are lost, whereas \method maintains a dynamic execution graph that models the internal effects of each block (tool calls with read $r$ or write $w$ effects, artifacts $v$, states st, failures F) and cross-block dependencies (consumes, equivalence, supersession), over which a declared intent selects continue, patch, replan, or answer.}
\label{fig:overview}
\vspace{-3mm}
\end{figure}


The concatenated representation produces three limitations that build on one another.
\label{sec:intro-limit}
At the representation layer, the history is a \emph{linear execution context}, where call provenance, artifact consumption, and state versioning have no unified form, so the agent rebuilds prior results and their relations from scratch.
At the analysis layer, feedback carries \emph{implicit execution progress}, because it never relates what a block intended to achieve to what it actually changed, and a block that raises no error can still re-fetch an existing result without advancing the task.
At the decision layer, standard output and errors give \emph{unstructured decision feedback} that describes only the current block, leaving the agent to infer from scattered history whether to proceed, fix the program, or change strategy, and it easily repeats a prior attempt or extends an unproductive path.

To address these limitations, we introduce \method, which keeps the programmatic expressiveness of PTC and maintains a task-level \emph{dynamic execution graph} built incrementally from real running traces.
As shown in Figure~\ref{fig:overview}, \method changes two structural aspects of PTC.
First, it models the internal logic of each block, so a code execution becomes a typed subgraph of tool calls with read or write effects, produced artifacts, state versions, and localized failures instead of an opaque output string.
Second, it models cross-block dependency as a multi-relation graph instead of a linear text stream, linking equivalent results, reused artifacts, and superseded states so that history becomes addressable.
Building on these two changes, \method runs three stages.
Graph-based execution modeling projects each trace into a shared typed graph and links cross-block dependencies.
Graph-based progress assessment binds a declared intent to the graph change it produces and separates new results from repeated ones.
Graph-based action selection turns the graph into a bounded summary over which the agent chooses to continue, patch a failed step, replan a dependency path, or answer.

We evaluate \method on seven tool-use benchmark settings with four backbone LLMs.
\method improves over both Direct Tool Calling and PTC on every benchmark, by $7.4$ points over PTC on average, and it reaches a given accuracy with $1.5$ to $3$ times fewer interaction rounds than Direct Tool Calling.
Its advantage also widens with task complexity, reaching $43$ points over Direct Tool Calling on the hardest tasks, and it transfers across closed and open backbones.
Our contributions are three-fold.
\begin{itemize}[leftmargin=1.2em, itemsep=-0.3ex, topsep=0ex]
    \item \textbf{Graph-Structured Tool Calling.} To the best of our knowledge, \method is the first framework that models programmatic tool calling as a task-level dynamic execution graph, capturing in-block tool effects and cross-block dependencies that a linear history discards.
    \item \textbf{Graph-Guided Execution.} We bind a declared intent to the graph change it produces and construct a bounded graph feedback, which lets the agent judge progress from graph effects and treat recovery as an addressable graph operation over continue, patch, replan, and answer.
    \item \textbf{Consistent Empirical Gains.} On seven settings and four backbones, \method surpasses Direct Tool Calling and PTC on the primary metric, widens its advantage on complex long-horizon tasks, and reduces interaction rounds at equal accuracy.
\end{itemize}
